\documentclass[aps,prd,twocolumn,superscriptaddress,nofootinbib]{revtex4-2}
\usepackage{amsmath,amssymb,amsthm}
\usepackage[colorlinks=true,linkcolor=blue,citecolor=blue,urlcolor=blue]{hyperref}

\newcommand{\SU}{\mathrm{SU}(2)}
\newcommand{\Tr}{\mathrm{Tr}}
\newtheorem{theorem}{Theorem}
\newtheorem{definition}{Definition}
\newtheorem{remark}{Remark}

\begin{document}

\title{A Certified Finite-Cutoff Program for the $\SU$ Theta-Graph
Transfer Operator: Interacting Positivity, Exact Seam-Integer Counting,
and a First Uniformity-in-the-Cutoff Theorem}

\author{Monty Dabas}
\affiliation{Celextrix Private Limited, New Delhi 110086, India}
\email{montydabas@celextrix.org}

\date{\today}

\begin{abstract}
We report a finite, machine-certified program studying the spectral gap of
the interacting transfer operator on the smallest non-trivial closed graph
with a plaquette --- the theta graph, two vertices joined by three edges
--- for $\SU$ lattice gauge theory. Every numerical statement below is an
exact two-sided rational (or exact-integer) enclosure produced by directed
interval arithmetic or exact linear algebra, never a floating-point
estimate; every claim carries a machine-checked certificate reproducible
from the accompanying repository. We prove: (i) the free reduced gap in
closed form via the exact $\SU$ character spectrum; (ii) a certified
two-sided interacting gap at every coupling $\kappa \geq 0$ on this graph,
via a Kill-Lemma-shaped decomposition into a block Schur estimate and a
positivity statement, both certified rather than assumed; (iii) that the
first-order crossing direction is an exact rank-one, graph-symmetry
--protected line, not an artifact of the compression used to compute it;
(iv) an exact eigenvalue-counting reformulation of the gap via a
Haynsworth-congruence (Birman--Schwinger-type) construction, giving exact
integer spectral counts rather than numerical brackets on part of the
coupling range; (v) the program's first statement that is uniform in the
lattice cutoff rather than proved at one fixed cutoff, obtained by
replacing the Wilson action with the heat-kernel action, whose exact
semigroup property removes discretization error identically; and (vi) that
the interacting gap is volume- and cutoff-uniform simultaneously on a
family of graphs built by vertex-gluing theta graphs with no shared
plaquettes, while the same route provably fails once plaquettes are
shared, localizing the residual volume obstruction exactly to
interaction overlap. We are explicit about what this program does not
establish: it does not touch the infinite-volume, continuum, or
Osterwalder--Schrader-reconstruction content of the Yang--Mills
existence-and-mass-gap problem, and none of the results here should be
read as bearing on that problem beyond providing an exactly verified
finite testbed and a catalogue of which certification routes do and do
not survive volume growth.
\end{abstract}

\maketitle

\section{Introduction}\label{sec:intro}

The Yang--Mills existence and mass-gap problem asks for a proof that a
non-trivial quantum Yang--Mills theory exists on $\mathbb{R}^4$ satisfying
the Wightman/Osterwalder--Schrader axioms, with a mass gap
$\Delta > 0$~\cite{JaffeWitten}. That statement concerns the continuum,
infinite-volume limit of the theory. The present paper does not attempt
it, and no theorem below should be misread as a step toward the
Clay-sense predicate.

What we do attempt is narrower and, we argue, still useful: to take the
smallest closed graph carrying a non-trivial plaquette --- the theta graph
$B_3$, two vertices joined by three edges, one independent holonomy after
gauge fixing, three faces --- and certify, with exact arithmetic and
machine-checked proofs, every spectral statement that can honestly be made
about its interacting transfer operator at finite lattice spacing. The
discipline we impose throughout is: \emph{no floating-point number is ever
part of a verdict}. Every enclosure below is either an exact rational
computation or a two-sided interval of provably bounded width obtained
from directed rounding, and every theorem is accompanied by an executable
certificate, a set of negative controls that must fail when the underlying
claim is deliberately falsified, and an explicit statement of what remains
open. This is closer in spirit to interval-certified numerics and to
formally verified mathematics than to a physics estimate, and we adopt it
because the target problem is famous partly for how easy it is to
overclaim.

The paper is organized as a linear certification chain, YM-1 through
YM-14, each step consuming only what the previous steps proved. We give
theorem statements, proof sketches sufficient for a reader to reconstruct
the argument, and pointers to the exact certificate in the accompanying
repository for full reproducibility. Section~\ref{sec:setup} fixes the
graph, the gauge-fixed carrier, and the character-expansion toolkit.
Section~\ref{sec:free} treats the free (non-interacting) reduced gap.
Section~\ref{sec:interacting} builds the interacting gap at finite
coupling on the theta graph in stages: a sandwich bound, the first-order
crossing direction, a symmetry-protection theorem, a certified two-sided
gap beyond the sandwich's reach, and an exact eigenvalue-counting
reformulation. Section~\ref{sec:capstone} states the capstone theorem: a
strictly positive gap at every coupling on the fixed graph, together with
the honest statement of what ``every coupling, fixed graph'' does not
give. Section~\ref{sec:uniform} gives the paper's central positive
result on the harder axis: a gap that is uniform in the lattice cutoff.
Section~\ref{sec:volume} extends this along the volume axis on a
restricted family and localizes exactly where the extension fails.
Section~\ref{sec:gates} audits three structural gates --- gauge
invariance, volume/infrared behaviour, and regulator universality --- and
reports two of the three as genuinely split between a closed and an open
half. Section~\ref{sec:open} states the remaining obligations without
euphemism. Section~\ref{sec:conclusion} concludes.

\section{Graph, carrier, and character toolkit}\label{sec:setup}

Let $B_3$ be the theta graph: two vertices $v_0,v_1$ joined by three edges
$e_A,e_B,e_C$ with $\SU$-valued holonomies $A,B,C$. Tree gauge-fixing on
any connected graph is exact and free of Gribov ambiguity: fixing a
spanning tree (here, edge $e_C$, so $C=\mathbb{1}$) leaves exactly
$b_1(B_3)=E-V+1=2$ independent holonomies, here $(A,B)$, with residual
gauge acting diagonally by conjugation. Gauge-invariant states are
therefore exactly $L^2(\SU\times\SU)^{\mathrm{Ad}}$, and this fact is
itself certified rather than assumed (Section~\ref{sec:gates}); every
capsule in this paper uses it as the carrier.

We work in the character basis. For $g\in\SU$, the irreducible
characters $\chi_j(g)$, $j\in\{0,\tfrac12,1,\tfrac32,\dots\}$, satisfy the
exact orthogonality and Clebsch--Gordan pairing relations of $\SU$; in
particular the triple integral
\begin{equation}
\int_{\SU}\chi_p(A)\,\chi_q(B)\,\chi_r(AB^{-1})\,dA\,dB
= \frac{\delta_{p=q=r}}{d_p},
\end{equation}
with $d_p=2p+1$, is exact and rational, and is the single identity that
governs both the intra- and inter-plaquette seam structure used
throughout. The Wilson single-link heat kernel has exact character
expansion $K_\beta(g)=\sum_j d_j\,\lambda_j(\beta)\,\chi_j(g)$ with
$\lambda_j(\beta)=I_{2j+1}(\beta)/I_1(\beta)$ (modified Bessel functions),
and the true heat-kernel (semigroup) action used from
Section~\ref{sec:uniform} on has exact expansion coefficients
$f_j(\kappa)=2I_{2j+1}(\kappa)/\kappa$ with the exact composition law
$K_a * K_b = K_{a+b}$ --- an identity the Wilson action does not share,
which is why the uniformity result of Section~\ref{sec:uniform} requires
the switch.

All numerical enclosures below use directed rational interval arithmetic
(no floating point in any verdict); Bessel-function values are enclosed
via their positive power series with a certified geometric tail bound, and
transcendental combinations (logs, exponentials) via certified two-sided
Taylor or continued-fraction enclosures with explicit remainder bounds.
Every enclosure states its width, and every stage includes at least one
deliberately falsified control (a wrong-recurrence Bessel tamper, a
coupling substitution, a sign flip) that must, and does, produce a
separated bracket.

\section{The free reduced gap}\label{sec:free}

\begin{theorem}[YM-1]\label{thm:ym1}
For the normalized single-link Wilson convolution operator with
$\beta=2$, the reduced spectral gap
\begin{equation}
\Delta_{\mathrm{red}}(a{=}1,\beta{=}2) = -\log\!\big(I_2(2)/I_1(2)\big)
\end{equation}
lies in the certified enclosure
$[0.836723306231588913050067801054\ldots]$, width $<10^{-30}$, an exact
rational interval.
\end{theorem}

Proof: $\lambda_j(\beta)=I_{2j+1}(\beta)/I_1(\beta)$ is exact by the
character expansion of $K_\beta$; $\lambda_{1/2}\in(0,1)$ strictly
(control C1), so the gap is strictly positive; a tampered substitution
$I_2\to I_3$ separates the bracket (control C2); a width budget is
enforced as a hard failure condition (control C3). This is the program's
base case and the only place where a bare single-link spectrum, rather
than a graph transfer operator, is the object of study.

\section{The interacting theta-graph gap at finite coupling}\label{sec:interacting}

The interacting transfer operator on $B_3$ is
$T_\kappa = M^{1/2}(K\times K)M^{1/2}$, where $K$ is the single-link
kernel of Section~\ref{sec:setup} and $M=M_\kappa$ is the plaquette
interaction weight $e^{(\kappa/2)\Tr U}$ on the third (closing) face.

\begin{theorem}[YM-2: sandwich bound]\label{thm:ym2}
For $0\le\kappa<\kappa_0=\Delta_{\mathrm{red}}/6$, the min-max sandwich
$m_-T_0\preceq T_\kappa\preceq m_+T_0$ holds with
$m_+/m_-=e^{6\kappa}$, using $|\Tr U|\le 2$. At $\kappa=1/8$ this
certifies a reduced gap $\ge 0.08672330623\ldots$.
\end{theorem}

\begin{theorem}[YM-3: first-order crossing direction]\label{thm:ym3}
The free top-excited eigenspace is exactly two-dimensional,
$\{\chi_{1/2}(A),\chi_{1/2}(B)\}$, certified separated from
$\chi_{1/2}^2,\chi_1,\chi_{3/2}$. The exact theta integral
$\int\chi_{1/2}(A)\chi_{1/2}(B)\chi_{1/2}(AB^{-1}) = 1/2$ splits this
doublet at first order into eigenvalues shifted by $\pm\lambda_{1/2}/4$;
the vacuum first derivative is exactly $0$. Hence
$r'(0)=\lambda_{1/2}/4\in[0.10828185668\ldots]$, and the crossing
direction is the single symmetric line
$\chi_{1/2}(A)+\chi_{1/2}(B)$ --- rank one, and invariant under the
graph's $A\leftrightarrow B$ symmetry.
\end{theorem}

\begin{theorem}[YM-4: symmetry protection]\label{thm:ym4}
Let $S:(A,B)\mapsto(B,A)$. Then $[S,T_\kappa]=0$ for \emph{every}
$\kappa$, because $\Tr(BA^{-1})=\Tr(AB^{-1})$ on $\SU$ (real characters).
Consequently the crossing line identified at first order in
Theorem~\ref{thm:ym3} cannot rotate out of the swap-even sector at any
coupling: this is a theorem, not a first-order accident. Certified
finite-$\kappa$ spectral data, obtained from the exact expansion
$e^{(\kappa/2)\Tr U}=\sum_j d_j f_j(\kappa)\chi_j$ with exact
Clebsch--Gordan ring arithmetic and a certified truncation remainder,
shows the swap-even branch rising and the swap-odd branch falling with
$\kappa$ on the sampled grid.
\end{theorem}

\begin{theorem}[YM-5: certified two-sided gap beyond the sandwich]\label{thm:ym5}
Using the exact doubling identity $m_\kappa^2=m_{2\kappa}$, a Weyl
bound $\|PMQ\|^2\le\lambda_{\max}(Pm_{2\kappa}P-(Pm_\kappa P)^2)$ and
$\|QSQ\|\le\lambda_{1/2}^2e^{3\kappa}$ certify a two-sided gap at
couplings beyond $\kappa_0$: at $\kappa=1/8$, ratio $\le0.5003$, gap
$\ge0.6924$; at $\kappa=1/4$ (past $\kappa_0$), ratio $\le0.5610$, gap
$\ge0.5781$; at $\kappa=3/10$, gap $\ge0.5332$; the route fails closed at
$\kappa=1$.
\end{theorem}

Theorem~\ref{thm:ym5} is, structurally, an unconditional realization of
the Schur-estimate-plus-positivity (``Kill Lemma'') skeleton of an
earlier, conjectural manuscript by the author that assumed the existence
of a vacuum state; no conjecture, claim, or continuum statement from that
manuscript is imported here (see the repository's \texttt{LINEAGE.md} for
the explicit non-import ruling). Only the block-Schur-plus-positivity
\emph{shape} of the argument survived, now realized with certified
exact-interval arithmetic on the graph carrier.

\begin{theorem}[YM-6/YM-7: exact seam-integer counting]\label{thm:ym67}
On the native five-dimensional carrier
$V_5=\{1,\chi_{1/2}(A),\chi_{1/2}(B),\chi_{1/2}(A)\chi_{1/2}(B),
\chi_{1/2}(AB^{-1})\}$ (all exact eigenvectors of $T_0$, exact Gram
matrix with a single off-diagonal overlap $1/2$), a Haynsworth
congruence --- the same finite-to-infinite compression device used
elsewhere in this program for a Birman--Schwinger-type reduction ---
turns the gap statement into an \emph{exact integer count}
$k_\Sigma(\kappa,\mu)=\#\{\mathrm{spec}(T_\kappa)>\mu\}$, computed via two
interval-LDL inertia evaluations of a rational pencil. Certified exact
counts: $k=1$ at $(\kappa,\mu)=(1/8,3/5),(1/4,3/5),(1/2,1)$, i.e.\
$\lambda_2\le\mu<\lambda_1$ \emph{exactly} --- a strictly stronger
statement than a numerical bracket. Enlarging the carrier to
$V_7=V_5\cup\{\chi_1(A),\chi_1(B)\}$ (Gram block-diagonal, all seven
vectors exact $T_0$-eigenvectors) lowers the complement top eigenvalue
from $\lambda_1$ to $\lambda_1\lambda_{1/2}$ and unlocks the previously
undecidable cell: exact count $k=1$ at $(7/10,13/10)$. On $V_5$ the cell
$(7/10,5/4)$ is honestly \textsc{refused}, bracket $[1,5]$: the carrier is
too small to decide it, and this is reported as a refusal rather than
resolved by widening tolerance.
\end{theorem}

\section{Capstone: a gap at every coupling on the fixed graph}\label{sec:capstone}

\begin{theorem}[YM-8]\label{thm:ym8}
The interacting theta-graph transfer operator has a strictly positive
spectral gap at \emph{every} coupling $\kappa\ge0$, on the fixed graph
$B_3$ at fixed lattice spacing.
\end{theorem}

Proof: a certified pointwise kernel floor
$k\ge e^{-3\kappa}\big[c_0^{-1}e^{-\beta}\big]^2>0$ (exact enclosures,
positive-series Bessel bounds) combined with the classical
Jentzsch/Krein-Rutman theorem~\cite{Jentzsch,KreinRutman} --- cited, not
rederived, per the citation discipline stated in
Section~\ref{sec:gates} --- gives a simple Perron eigenvalue, hence a
strict gap, at every $\kappa$. The window $\kappa\in[0,7/10]$ additionally
carries the exact counts and enclosures of Theorems~\ref{thm:ym2}
--\ref{thm:ym67}.

\begin{remark}
We state as plainly as possible what Theorem~\ref{thm:ym8} does
\emph{not} say. A positive gap at every fixed coupling on a fixed,
finite graph is compatible with the gap closing as the graph is refined
or as the lattice spacing is taken to zero. Nothing here is uniform in
the cutoff or the volume; that is exactly the content that
Sections~\ref{sec:uniform}--\ref{sec:volume} address on restricted
families, and exactly the content that remains open in general
(Section~\ref{sec:open}). A later result in this program
(Theorem~\ref{thm:ym12}) further demotes the Jentzsch/Krein-Rutman
anchor used here to a \emph{simplicity-only} role, once positivity
itself is re-derived from a manifest square; that correction is recorded
in the governance ledger of the repository and does not alter
Theorem~\ref{thm:ym8}'s conclusion.
\end{remark}

\section{A first uniformity-in-the-cutoff theorem}\label{sec:uniform}

Every result so far is at one fixed lattice spacing $a$. We now obtain
the program's first result that holds \emph{for every} $a>0$
simultaneously, by replacing the Wilson action with the heat-kernel
action $K_a(g)=\sum_j d_j\,e^{-aC_j}\,\chi_j(g)$, $C_j=j(j+1)$ the
Casimir. This action satisfies the exact semigroup law
$K_a * K_b = K_{a+b}$: composing $n$ links of spacing $a/n$ reproduces
one link of spacing $a$ \emph{exactly}, with no discretization error.
The Wilson action has no such exact composition law, which is why it
cannot support the argument below.

\begin{theorem}[YM-9]\label{thm:ym9}
The free reduced gap under the heat-kernel action is
$\Delta_{\mathrm{free}}(a) = C_{1/2} = 3/4$ exactly, for every $a>0$ ---
the transcendental $a$-dependence cancels identically. Along the
declared linear trajectory $\kappa(a)=\theta a$, the interacting gap
satisfies $\Delta(a,\kappa(a))\ge 3/4-6\theta$ for every $a>0$; at
$\theta=1/16$ this gives a cutoff-independent gap of exactly $3/8$,
failing closed at $\theta\ge1/8$. The seam count
$k(a,e^{-as})$ is exactly $a$-independent (a combinatorial fact about
which Casimir levels lie below a threshold), so the exact-counting
machinery of Theorem~\ref{thm:ym67} transports across the entire
refinement family at once, not just at one spacing. By contrast, at
fixed Wilson coupling $\beta$ the reduced gap grows like $1/a$ and
diverges as $a\to0$ --- there is no cutoff-independent gap on this action
without a compensating running coupling.
\end{theorem}

We state the honest remainder of Theorem~\ref{thm:ym9} without
softening. The trajectory $\kappa(a)=\theta a$ is \emph{declared}, not
derived: the physically motivated trajectory is fixed by asymptotic
freedom, $\beta\sim\log(1/a)$, which is not modelled here. The graph is
fixed, so this is a statement about refining the same graph's lattice
spacing, not about infinite-volume growth. The heat-kernel action is
itself a choice among a regulator class (Section~\ref{sec:gates}
addresses which parts of the program are regulator-independent). And
$\Delta$ throughout is a reduced-graph transfer-operator gap, not a
reconstructed physical mass gap in any continuum sense. What has moved is
narrow but real: one structural gate --- uniformity in the cutoff ---
goes from wholly untouched to touched, with an exact theorem, on a toy
carrier.

\section{Volume: where uniformity extends and where it provably fails}\label{sec:volume}

\begin{theorem}[YM-12, volume half]\label{thm:ym12}
Let $L_m$ be the graph formed by vertex-gluing $m$ copies of the theta
graph with \emph{no shared plaquettes}. The transfer operator
tensor-factorizes exactly across the $m$ copies, so the ratio
$\lambda_2/\lambda_1$ is $m$-independent, and combined with
Theorem~\ref{thm:ym9}, $\Delta(L_m,a,a/16)\ge3/8$ for \emph{every} $m$
and \emph{every} $a$ --- a gap uniform in volume and cutoff
simultaneously, on this family.
\end{theorem}

This closes the interacting half of the volume gate audited in
Theorem~\ref{thm:ym11} below, but only on the no-shared-face family.
Theorem~\ref{thm:ym11} shows the same sandwich-style route fails, exactly
and quantifiably, once faces are shared; we record the two extremes as
the boundaries of a genuine open problem, not as a resolved one.

The same paper (YM-12) also re-derives the positivity underlying
Theorems~\ref{thm:ym5}--\ref{thm:ym8} from a manifest square,
$T_\kappa(a)=S^*S$ with $S=[K_{a/2}\otimes K_{a/2}]\,m_{\kappa/2}$ ---
both factors being the program's own halved-parameter objects from
Theorem~\ref{thm:ym9}'s semigroup property and Theorem~\ref{thm:ym5}'s
doubling identity. This demotes the classical Jentzsch/Krein-Rutman
citation used in Theorem~\ref{thm:ym8} to a role of certifying only
\emph{simplicity} of the top eigenvalue, with positivity itself now
square-sourced from within the program's own objects rather than
imported.

\section{Three structural gates, audited}\label{sec:gates}

We close the finite program with an audit of three structural questions
on the toy-carrier family $B_n$ (the ``$n$-banana'' graph: two vertices,
$n$ edges, $b_1=n-1$ holonomies, $F(n)=n(n-1)/2$ plaquettes; $B_3$ is the
theta graph used throughout).

\begin{theorem}[YM-11, Gate 1 --- gauge: closed]\label{thm:ym11a}
Tree gauge-fixing on a connected graph is exact, with no Gribov
obstruction: gauge-invariant states are exactly
$L^2(G^{b_1})^{\mathrm{Ad}}$. For $B_3$ this certifies, rather than
assumes, the carrier used in every capsule above.
\end{theorem}

\begin{theorem}[YM-11, Gate 2 --- volume/infrared: split]\label{thm:ym11}
The free reduced gap $C_{1/2}=3/4$ is volume- and cutoff-uniform for
every $n$ and $a$: closed. The interacting sandwich bound gives
$\Delta\ge C_{1/2}-2F(n)\theta$, which degrades \emph{quadratically} in
$n$: at $\theta=1/16$ the certified bounds are $3/8$ ($n{=}3$), exactly
$0$ ($n{=}4$), and $-1/2$ ($n{=}5$). The sandwich/envelope route is
therefore proven \emph{insufficient} for volume uniformity beyond
Theorem~\ref{thm:ym12}'s no-shared-face restriction, with the critical
volume computed exactly rather than estimated.
\end{theorem}

\begin{theorem}[YM-11, Gate 3 --- regulator universality: split]\label{thm:ym11b}
For the whole regulator class $K=\sum_j d_j c_j\chi_j$ with $c_j>0$
(Wilson, heat-kernel, and others), the \emph{counting} layer --- content
grading, the exact multiplicity law $m(j_1,j_2)=\min(2j_1,2j_2)+1$ on the
Ad-invariant sector, carrier dimensions, and the source-restriction
``blindness'' ledger of Theorem~\ref{thm:ym10} below --- depends only on
$\SU$ representation theory and is therefore regulator-independent:
closed. The \emph{metric} layer, i.e.\ the numerical value of the gap,
genuinely differs between regulators (witnessed directly by the Wilson
vs.\ heat-kernel contrast in Theorem~\ref{thm:ym9}) and remains open.
\end{theorem}

\begin{theorem}[YM-10, source-restriction ledger]\label{thm:ym10}
Every compressed carrier $V$ used above is, in the sense of the
companion source-restriction framework, a restriction of the full target
space; the relevant question is not tightness of the resulting bound but
whether $V$ is \emph{blind} to the target. The exact multiplicity law
above independently reproduces $\dim V_5=5$ and $\dim V_7=7$ and
re-derives the two-dimensionality of the $(\tfrac12,\tfrac12)$ sector
used in Theorem~\ref{thm:ym3} from representation theory rather than from
the Gram matrix. The exact probe count $b(V,s)=\mathrm{rank}(\Pi_s\mid
\ker E)$ shows $b(V_5,s)=0$ for $s\le2$ --- the carrier used in
Theorem~\ref{thm:ym67} is exactly blindness-free within its own certified
threshold window, which is why that choice of carrier was correct and not
merely convenient. At $s=3$, $V_5$ requires six probes and $V_7$ four;
blindness is nonincreasing along $V_5\subset V_7\subset V_9$ with
nonnegative stage increments; and the whole ledger, depending only on
content and Casimir level, is exactly $a$-independent, so one ledger
covers the entire refinement family of Theorem~\ref{thm:ym9}. This ledger
is exact for the free target only; on interacting faces the content
grading mixes, so $b(V,s)$ is a lower bound there, and closing that gap
is an explicit named obligation of the program, not a silent one.
\end{theorem}

We record one self-correction as part of the certification discipline
itself: an earlier draft of Theorem~\ref{thm:ym9}'s uniform seam count
(T4) counted spectral \emph{contents}, while the counting statement of
Theorem~\ref{thm:ym67} counts \emph{eigenvalues with multiplicity}; at
the relevant sector these differ ($4$ versus $5$). Both quantities are
independently $a$-independent, so the uniformity conclusion of
Theorem~\ref{thm:ym9} is unaffected, but the arithmetic was wrong and is
corrected here. We regard catching and publishing this kind of error,
rather than avoiding it, as part of what the certification discipline is
for.

\section{Interaction overlap between plaquettes}\label{sec:overlap}

\begin{theorem}[YM-14]\label{thm:ym14}
For a bridged pair of theta-graph blocks coupled by a single shared
face, $e^{(\kappa/2)\Tr(A_1A_2^{-1})}$ --- the minimal case of genuine
plaquette overlap not covered by Theorem~\ref{thm:ym12}'s no-sharing
family --- the inter-block overlap seam is exactly rank-one and
transported by the same graph symmetry and the same theta integral
($1/2$) that governs the intra-block seams of Theorem~\ref{thm:ym3}: one
identity governs both. On the orthonormal carrier
$W_5=\{1,\chi_{1/2}(A_1),\chi_{1/2}(A_2),\chi_{1/2}(B_1),\chi_{1/2}(B_2)\}$,
exact counts $k=1$ are certified at
$(\kappa,\mu)=(1/8,3/5),(1/4,7/10),(1/2,9/10)$; at $\kappa=1$ the carrier
certifies count $0$ at $\mu=3/2$ (the compressed vacuum value is below
$\mu$), which we report as a carrier limitation rather than resolve by
enlarging tolerance. The first-order overlap price is exactly
$\kappa\lambda/4$, carried by a single swap-even line, in direct analogy
with the single-block crossing direction of Theorem~\ref{thm:ym3}.
\end{theorem}

This opens, but does not close, the object named by
Theorem~\ref{thm:ym11}'s Gate 2 as the real residual obstruction:
Theorem~\ref{thm:ym12} is uniform on no-shared-face volume growth,
Theorem~\ref{thm:ym11} shows the sandwich route fails on complete
sharing, and a physically realistic bounded-degree lattice sits strictly
between these two certified extremes. The next object named by the
program is a one-dimensional block-transfer dock for a chain of
bounded, but nonzero, overlap; it is not attempted here.

\section{What remains open}\label{sec:open}

We list the open content without euphemism, since the value of this
program depends on the list being accurate.

\begin{itemize}
\item \textbf{Volume, general lattice.} Theorem~\ref{thm:ym12} is uniform
only on the no-shared-face family; Theorem~\ref{thm:ym11} shows the
sandwich route fails, with an exactly computed critical volume, once
faces are shared; the bounded-overlap case of physical interest
(Theorem~\ref{thm:ym14}) is opened but not closed.
\item \textbf{Regulator universality, metric layer.} Only the
combinatorial/counting layer of universality (Theorem~\ref{thm:ym11b})
is closed; the numerical gap value differs by regulator and the
asymptotically-free running-coupling trajectory needed to connect the
heat-kernel result to the physical continuum limit is declared, not
derived (Theorem~\ref{thm:ym9}).
\item \textbf{Osterwalder--Schrader reconstruction.} No statement here
constructs, or claims to construct, a continuum Wightman theory from
these finite-graph transfer operators.
\item \textbf{Non-triviality and the Clay predicate.} None of the above
bears on the existence-and-mass-gap predicate of the Clay Millennium
Problem as stated in~\cite{JaffeWitten}; that problem concerns the
continuum, infinite-volume limit, which this program has not
addressed and does not claim to have addressed.
\end{itemize}

\section{Conclusion}\label{sec:conclusion}

We have given, on the smallest graph carrying a Yang--Mills plaquette, a
fully exact-arithmetic, machine-certified account of the interacting
transfer operator's spectral gap: a strictly positive gap at every
coupling on the fixed graph (Theorem~\ref{thm:ym8}), an exact
eigenvalue-counting reformulation on part of the coupling range
(Theorem~\ref{thm:ym67}), the program's first cutoff-uniform theorem
(Theorem~\ref{thm:ym9}), and a volume-uniform extension on a restricted
family with the extension's exact failure point identified on the
complementary family (Theorems~\ref{thm:ym12}, \ref{thm:ym11}). Each
result is accompanied by negative controls that must, and do, fail when
the claim is deliberately falsified, and each is reproducible from the
certificates and scripts in the accompanying repository. The value we
claim for this program is not proximity to the continuum problem, which
remains entirely open, but a fully verified, error-caught, and
error-corrected finite testbed --- together with a precise catalogue,
Section~\ref{sec:open}, of exactly which structural gates a genuine
continuum proof would still need to close.

\section*{Data and code availability}

All certificates, verification scripts, and test suites are available in
the public repository \texttt{Parveen117/Publications}, directory
\texttt{papers/yang-mills-certified-benchmark/}. Each theorem above
corresponds to a numbered script (\texttt{ym1\_certified\_gap.py}
through \texttt{ym14\_overlap\_dock.py}) whose execution regenerates and
pins the stated certificate; the repository's continuous-integration
workflow reruns and pin-diffs every certificate on each change.

\begin{thebibliography}{9}

\bibitem{JaffeWitten}
A.~Jaffe and E.~Witten, ``Quantum Yang--Mills Theory,'' Clay Mathematics
Institute Millennium Prize Problem description (2000).

\bibitem{Jentzsch}
R.~Jentzsch, ``Über Integralgleichungen mit positivem Kern,''
J.\ Reine Angew.\ Math.\ \textbf{141}, 235 (1912).

\bibitem{KreinRutman}
M.~G.~Krein and M.~A.~Rutman, ``Linear operators leaving invariant a
cone in a Banach space,'' Uspehi Mat.\ Nauk \textbf{3}, 3 (1948);
English translation, Amer.\ Math.\ Soc.\ Translation \textbf{26} (1950).

\bibitem{ReedSimonIV}
M.~Reed and B.~Simon, \emph{Methods of Modern Mathematical Physics IV:
Analysis of Operators} (Academic Press, 1978), Theorem XIII.43.

\bibitem{Haynsworth}
E.~V.~Haynsworth, ``Determination of the inertia of a partitioned
Hermitian matrix,'' Linear Algebra Appl.\ \textbf{1}, 73 (1968).

\bibitem{Wilson}
K.~G.~Wilson, ``Confinement of quarks,'' Phys.\ Rev.\ D \textbf{10},
2445 (1974).

\end{thebibliography}

\end{document}
